% !TEX root = ../../main.tex
% 条件溯源概念流程图（第二章，简化版主干流程）
\begin{figure}[H]
\centering
\begin{tikzpicture}[
  node distance = 0.55cm and 1.6cm,
  proc/.style   = {rectangle, draw, rounded corners=3pt, minimum width=3.0cm, minimum height=0.65cm,
                   font=\small, fill=blue!8, draw=blue!50!black},
  decision/.style = {diamond, draw, aspect=2.2, minimum width=3.2cm, minimum height=0.8cm,
                     font=\small, align=center, fill=orange!15, draw=orange!60!black, inner sep=2pt},
  term/.style   = {rectangle, draw, rounded corners=8pt, minimum width=3.0cm, minimum height=0.65cm,
                   font=\small, fill=green!12, draw=green!50!black},
  reject/.style = {rectangle, draw, rounded corners=8pt, minimum width=2.2cm, minimum height=0.62cm,
                   font=\small, fill=red!12, draw=red!50!black},
  arr/.style    = {-Stealth, thick},
  lbl/.style    = {font=\small},
]

\node[proc]   (start)  {投诉申请（orderID + 原因）};
\node[decision, below=0.8cm of start] (d1) {完整性\\校验通过？};
\node[reject,  right=2.2cm of d1]  (r1) {拒绝溯源（日志可疑）};
\node[decision, below=0.9cm of d1] (d2) {关键词\\命中？};
\node[reject,  right=2.2cm of d2]  (r2) {拒绝溯源（无据可查）};
\node[proc,    below=0.9cm of d2]  (pid) {PID 溯源};
\node[term,    below=0.7cm of pid] (uid) {获得 userID（真实身份）};
\node[proc,    below=0.6cm of uid] (log) {写入审计日志};

\draw[arr] (start) -- (d1);
\draw[arr] (d1) -- node[lbl, left]{是} (d2);
\draw[arr] (d1) -- node[lbl, above]{否} (r1);
\draw[arr] (d2) -- node[lbl, left]{是} (pid);
\draw[arr] (d2) -- node[lbl, above]{否} (r2);
\draw[arr] (pid) -- (uid);
\draw[arr] (uid) -- (log);

\end{tikzpicture}
\caption{条件溯源概念流程（预备知识层）}
\label{fig:trace_concept}
\end{figure}
